% =====================================================================
\section{Special function registers and flags}
\label{sec:sfr}
% =====================================================================

Control and peripheral registers reside in the 8-bit SFR address space, which
is reachable through the SFR addressing mode of every two-operand instruction
(e.g.\ \texttt{OR R1, [SFR:I2CCON]}). Peripherals are therefore operated with
the same instructions as memory.

\subsection{SFR map}

{\small
\begin{xltabular}{\linewidth}{@{}l l l l X@{}}
  \toprule
  \textbf{Address} & \textbf{Name} & \textbf{Access} & \textbf{Reset} & \textbf{Description} \\
  \midrule
  \endfirsthead
  \toprule
  \textbf{Address} & \textbf{Name} & \textbf{Access} & \textbf{Reset} & \textbf{Description} \\
  \midrule
  \endhead
  \bottomrule
  \endfoot
  \hex{00} & SR      & R/W   & \hex{00} & Status register (flags, IE), see \ref{sec:sr} \\
  \hex{01} & SPL     & R/W   & see \reqref{INT-011} & Stack pointer, low byte \\
  \hex{02} & SPH     & R/W   & see \reqref{INT-011} & Stack pointer, high byte \\
  \hex{03} & CWP     & R/W   & \hex{00} & Bits 3:0: current register window \\
  \hex{04} & WSTART  & R/W   & \hex{00} & Bits 3:0: oldest window held in the register file \\
  \hex{05} & PRIDX   & R/W   & \hex{00} & Bits 5:0: physical register index \\
  \hex{06} & PRDAT   & R/W   & --       & Read/write of physical register P[PRIDX] \\
  \hex{08} & INTEN   & R/W   & \hex{00} & Maskable interrupt enable (bit 0 INT0, bit 1 \IIC) \\
  \hex{09} & INTFLG  & R/W1C & \hex{00} & Pending interrupt flags; writing 1 clears \\
  \hex{10} & I2CCON  & R/W   & \hex{00} & \IIC{} controller control, see Section~\ref{sec:i2c} \\
  \hex{11} & I2CSTA  & R/W1C & \hex{00} & \IIC{} controller status \\
  \hex{12} & I2CADR  & R/W   & \hex{00} & Bits 6:0: 7-bit slave address \\
  \hex{13} & I2CTXD  & W     & \hex{00} & Byte to transmit \\
  \hex{14} & I2CRXD  & R     & \hex{00} & Last received byte \\
  \hex{15} & I2CBRG  & R/W   & \hex{FF} & SCL clock divider \\
  \multicolumn{4}{@{}l}{\hex{07}, \hex{0A}--\hex{0F}, \hex{16}--\hex{FF}} & Reserved; reads return \hex{00}, writes are ignored \\
\end{xltabular}
}

\subsection{Status register SR}
\label{sec:sr}

\begin{center}
\small
\begin{tabularx}{0.85\linewidth}{@{}l l X@{}}
  \toprule
  \textbf{Bit} & \textbf{Name} & \textbf{Meaning} \\
  \midrule
  0   & C  & Carry / borrow; for \IIC{} instructions: transaction error \\
  1   & Z  & Result is zero \\
  2   & N  & Bit 7 of the result (sign) \\
  3   & V  & Signed (two's complement) overflow \\
  4--6 & -- & Reserved, read as 0 \\
  7   & IE & Global enable for maskable interrupts \\
  \bottomrule
\end{tabularx}
\end{center}

\subsection{Requirements}

\begin{reqtable}
\req{SFR-001}{The SFR address space \MUST{} contain 256 addresses (8 bits) and
  \MUST{} be separate from data memory.}
\req{SFR-002}{The SFRs \MUST{} be implemented as listed in the SFR map,
  including their reset values.}
\req{SFR-003}{A read from a reserved address \MUST{} return \hex{00}; a write
  to it \MUST{} be ignored without raising an exception.}
\req{SFR-004}{Bits marked R/W1C \MUST{} be cleared by writing 1; writing 0
  \MUST{} leave them unchanged.}
\req{SFR-005}{SPL/SPH \SHOULD{} be written in the order SPL, SPH; the
  implementation \MAY{} block interrupts for one instruction after a write to
  SPL to avoid an inconsistent SP.}
\req{SFR-006}{ALU instructions \MUST{} update the flags as listed in the
  ``Flags'' column of the instruction tables (Section~\ref{sec:isa}); all other
  instructions \MUSTNOT{} change the flags unless stated otherwise.}
\end{reqtable}
